% ============================================================
%  AI Tools Usage Declaration
%  Compile: pdflatex ai_declaration.tex
% ============================================================
\documentclass[10pt,a4paper]{article}

\usepackage[margin=1.8cm]{geometry}
\usepackage{booktabs}
\usepackage{array}
\usepackage{microtype}
\usepackage{parskip}
\usepackage{setspace}
\usepackage[hidelinks]{hyperref}

\setstretch{1.15}
\setlength{\parindent}{0pt}
\setlength{\parskip}{0.5em}

\begin{document}

% ── Header ───────────────────────────────────────────────────
\begin{center}
{\large\textbf{AI Tools Usage Declaration}}\\[0.3em]
{\small The Climate Standoff: A Dynamic Game Model of Strategic Tipping
in Global Climate Coordination}\\[0.2em]
{\small David Clements \quad|\quad Academic Year 2026}
\end{center}

\smallskip
\hrule
\smallskip

% ── Narrative ────────────────────────────────────────────────
AI tools were used throughout this dissertation as part of an iterative research
support workflow. Their primary role was to assist with debugging computational
implementation, stress-testing methodological choices, and accelerating
familiarisation with new technical tools and modelling concepts. All model design
decisions, calibration choices, and interpretation of results remained my own, and
code suggestions were reviewed, adapted, and integrated only where they aligned with
the intended equilibrium structure and economic logic of the project.

AI was also used editorially to improve clarity and presentation, particularly by
identifying areas of unnecessary repetition, excessive verbosity, or places where
further explanation would strengthen the argument. In later drafting stages, it was
used to simulate critical examiner-style feedback, helping expose weak justifications,
unclear assumptions, and potential methodological vulnerabilities prior to submission.

\textbf{AI tools used:} Claude Opus, Claude Sonnet, GPT-4, GPT-4o, Perplexity.

\smallskip

% ── Table ────────────────────────────────────────────────────
\begin{table}[h]
\small
\renewcommand{\arraystretch}{1.35}
\begin{tabular}{>{\bfseries}p{3.6cm} p{12.2cm}}
\toprule
Usage & Representative Prompts \\
\midrule

Literature discovery
  & ``I'm working on a dynamic game model of climate coordination with heterogeneous blocs and QRE.
    What are some must-read papers I might have missed?'' /
    ``Find recent applications of SMM in environmental economics'' /
    ``What papers cite Barrett (1994) and Harstad (2016) together?'' \\[4pt]

Learning \& teaching
  & ``Teach me quantal response equilibrium'' /
    ``Explain the difference between Nash and QRE in plain English'' /
    ``What is spectral radius and why does it matter for equilibrium uniqueness?'' /
    ``Have I correctly understood and cited this paper?'' \\[4pt]

Writing feedback
  & ``Does this read okay or is it too verbose?'' /
    ``Mark this dissertation and tell me what the worst flaws are'' /
    ``Where did I not go into enough detail?'' /
    ``Suggest a clearer way to phrase this policy implication'' \\[4pt]

Code generation \& correction
  & ``Write a function that computes the spectral radius of the logit response Jacobian'' /
    ``Check this fixed-point iteration for off-by-one errors'' /
    ``Rewrite this loop to be vectorised and faster'' \\[4pt]

Methodology review \& debugging
  & ``Do you see any glaring methodological flaws here?'' /
    ``Is my identification strategy for the SMM valid?'' /
    ``Why is the fixed-point iteration not converging at states with 4 active players?'' \\[4pt]

Proofreading \& grammar
  & ``Proofread this paragraph'' /
    ``Fix the referencing format here'' \\

\bottomrule
\end{tabular}
\end{table}

\end{document}
